\begin{pastexam}[label={px:ntu-108-ee-dsb-14}]{臺大電機所 108 學年度（B）第 14 題}
（複選；每題 10 分，每個選項單獨計分，答對 2 分，答錯倒扣 2 分）

An algorithm is said to be \textit{insensitive to input} if its running time is independent of the
initial state of its input. For sorting algorithms, that means that the run time depends only on
the number of elements to be sorted, not on their initial arrangement. Which sorting algorithms
are insensitive to input?
\begin{enumerate}[label=(\Alph*).]
\item Bubble sort.
\item Selection sort.
\item Merge sort.
\item Quick sort.
\item Radix sort.
\end{enumerate}
\end{pastexam}

\begin{pastsoln}
答案是 (B)(C)(E)。

先說判斷標準：「執行時間只和 $n$ 有關」要看漸近的時間，也就是最好和最壞情況是不是同一個
$\Theta$。 如果連確切的步數都要一模一樣，五種演算法幾乎都會因為比較或搬動的次數有小差異而
不算，題目就沒有意義了。 各演算法的最好、最壞時間見資結書第~14.7~節的比較表。

\textbf{(A) 錯}。 氣泡排序一般會加「這一趟沒有交換就結束」的判斷（資結書虛擬碼~14.1）：
輸入已經排好時只跑一趟，$\Theta(n)$；反序時跑滿 $n - 1$ 趟，$\Theta(n^2)$。 時間取決於輸入的
排列。 即使不提早結束，交換次數也等於逆序數，從 $0$ 到 $n(n-1)/2$ 不等。

\textbf{(B) 對}。 選擇排序第 $i$ 趟一定要掃過剩下的 $n - i$ 個元素找最小值，不管輸入怎麼排，
比較次數都是 $\sum_{i=1}^{n-1} (n - i) = n(n-1)/2$，交換最多 $n - 1$ 次，時間永遠是
$\Theta(n^2)$（資結書第~14.2~節）。

\textbf{(C) 對}。 合併排序不管輸入怎麼排，都對半切、遞迴、合併，遞迴式固定是
$T(n) = 2T(n/2) + \Theta(n)$，時間永遠是 $\Theta(n \lg n)$（資結書第~14.4~節）。 合併時的比較
次數會隨輸入略有不同（一邊先用完就不用再比），但每一層的搬動都是 $n$ 個，不影響 $\Theta$。

\textbf{(D) 錯}。 快速排序取決於基準把陣列切得平不平均：每次都切成一半時是 $\Theta(n \lg n)$；
取第一個當基準而輸入已經排好時，每次只切掉一個，$\Theta(n^2)$（資結書第~14.5~節）。

\textbf{(E) 對}。 基數排序對 $d$ 個位數各做一趟穩定的計數排序，每趟 $\Theta(n + k)$，總共
$\Theta(d(n + k))$（資結書第~15.2~節）。 每一趟都把所有元素分配、收集一次，做的事和元素原本的
順序無關。 （$d$ 由鍵的位數決定，和排列無關。）
\end{pastsoln}
